% Copyright (c) 2026 by C. Klukas.
% Licensed under the MIT License. See LICENSE for details.

\Needspace{8\baselineskip}
\subsection{State capture, lifecycle ordering and recovery}\label{sec:lifecycle}

Migration, cloning and deactivation take effect after the current handler
returns. A handler records its request; the host then captures requested
clones, runs the original instance's end-of-call hook, and finally
captures or retires the original. A clone therefore sees state from
before the original's departure or disposal hook changes it. No image is
taken while a handler or imported function is running. The implementation
is in \sourcefile{cpp/host/src/runtime.cpp}{\code{Runtime::Impl::finish\_call}}.

A paglet can work indefinitely, but cannot process another message or
capture an image during an uninterrupted handler. The handler watchdog is
off by default; operators may enable it for shared hosts. Long operations
therefore normally run in background services or compute jobs, with
continuations delivered in later handlers. This keeps status, cancellation
and state capture available while work proceeds. A handler budget, a
caller's reply timeout and an overall task deadline serve different
purposes; none is a general agent lifetime limit.

A runtime execution error (a \emph{trap}) or an overrun of an enabled time
budget in a handler or lifecycle hook ends the paglet as failed. If the dispatching
hook fails, the runtime does not publish the departure. Disposal and
deactivation are refused while an arrival or move is undecided, so
administrative operations cannot bypass the transfer protocol. Negative
return codes from ABI events are logged; returning such a code does not
veto a lifecycle transition.

The runtime delivers the initial lifecycle event before ordinary
messages: \code{created} for a new instance, \code{cloned} for a copy, and
\code{arrived} for a migrated instance. Migration does not also produce
\code{activated}. Request replies and host failure reports are delivered
later through the ordinary mailbox, the instance's message queue.
A request timeout calls its continuation with an error; it does not stop
a remote handler that is already running.

When persistent storage and restart are enabled, completed calls can
save a \emph{checkpoint}: the memory image and associated host record
needed for recovery. After a restart, activation and the optional
\code{recovered} event arrive before restored request timeouts and
\code{gone} replies for pending jobs. The ordinary mailbox is not saved
as a durable message queue. Effects after the last checkpoint may repeat
when execution resumes. WASI reactor initialization runs when the instance
is created, before host imports are attached; restoring an image then
replaces the memory produced by initialization. Constructors therefore do
not serve as a recovery protocol. Saving a checkpoint also does not make
changes to an external service part of the same atomic transaction.
